\chapter{Vergleichen und Ersetzen von Teilmengen}
\label{b13guide:ch:comparison}

\section{Vergleich von Infima und Suprema}

\begin{theorem}[Monotonie bezüglich der Teilmenge]
\label{b13guide:thm:monotone}
Für $S\subseteq T\subseteq P$ gilt bei Existenz der jeweils
beiden Werte
\[
 \sup S\leq\sup T,\qquad \inf T\leq\inf S.
\]
Die Existenz eines Supremums oder Infimums von $T$ garantiert
im allgemeinen nicht die Existenz des entsprechenden Wertes von $S$.
\end{theorem}


\begin{theorem}[Exakte Vergleichskriterien durch Schrankenmengen]
\label{b13guide:thm:exact-comparison}
Für $S,T\subseteq P$ gilt bei Existenz der jeweils beiden Werte
\[
\begin{aligned}
 \sup S\leq\sup T&\ \Longleftrightarrow\ \Hi(T)\subseteq\Hi(S),\\
 \inf S\leq\inf T&\ \Longleftrightarrow\ \Lo(S)\subseteq\Lo(T).
\end{aligned}
\]
Eine Inklusion zwischen $S$ und $T$ selbst ist dafür nicht erforderlich.
\end{theorem}

\begin{theorem}[Vergleich durch elementweise Überdeckung]
\label{b13guide:thm:domination}
Seien $S,T\subseteq P$. Gilt
\[
 \forall x\in S\;\exists y\in T\;(x\leq y),
\]
so ist $\Hi(T)\subseteq\Hi(S)$. Sind beide Suprema vorhanden,
so folgt $\sup S\leq\sup T$.
Gilt stattdessen
\[
 \forall x\in S\;\exists y\in T\;(y\leq x),
\]
so ist $\Lo(T)\subseteq\Lo(S)$. Sind beide Infima vorhanden,
so folgt $\inf T\leq\inf S$.

Gilt die obere Überdeckungsbedingung in beiden Richtungen, so ist
$\Hi(S)=\Hi(T)$; ein Supremum existiert genau dann, wenn das
andere existiert, und die Werte sind dann gleich. Gilt die untere
Überdeckungsbedingung in beiden Richtungen, so ist
$\Lo(S)=\Lo(T)$; dasselbe gilt für die Infima.

\end{theorem}

\begin{theorem}[Getrennte Mengen und gemeinsame Elemente]
\label{b13guide:thm:separation}
Seien $S,T\subseteq P$. Sind $\sup S$ und $\inf T$ vorhanden, so gilt
\[
 \sup S\leq\inf T
 \quad\Longleftrightarrow\quad
 \forall x\in S\;\forall y\in T\;(x\leq y).
\]
Ist $S\cap T\ne\varnothing$, und existieren $\inf S$ sowie $\sup T$,
so gilt $\inf S\leq\sup T$.
\end{theorem}

\section{Gleiche Schranken und unveränderte Werte}

\begin{theorem}[Gleiche Schrankenmengen]
\label{b13guide:thm:same-bounds}
Aus $\Hi(S)=\Hi(T)$ folgt: $\sup S$ existiert genau dann, wenn
$\sup T$ existiert; in diesem Fall sind beide gleich.
Entsprechend liefert $\Lo(S)=\Lo(T)$ äquivalente Existenz und
Gleichheit der Infima.
Sind beide Suprema bereits vorhanden, gilt auch die Umkehrung:
\[
 \sup S=\sup T\quad\Longleftrightarrow\quad\Hi(S)=\Hi(T).
\]
Für zwei vorhandene Infima gilt entsprechend
\[
 \inf S=\inf T\quad\Longleftrightarrow\quad\Lo(S)=\Lo(T).
\]
\end{theorem}

\begin{definition}[Kofinale und koinitiale Teilmenge]
\label{b13guide:def:cofinal}
Eine Teilmenge $S\subseteq T\subseteq P$ heißt \emph{kofinal in $T$},
wenn
\[
 \forall t\in T\;\exists s\in S\;(t\leq s).
\]
Sie heißt \emph{koinitial in $T$}, wenn
\[
 \forall t\in T\;\exists s\in S\;(s\leq t).
\]
Diese Begriffe beziehen sich hier ausschließlich auf das Ersetzen
einer Teilmenge bei der Bestimmung ihrer Schranken.
\end{definition}

\begin{theorem}[Kofinale und koinitiale Reduktion]
\label{b13guide:thm:cofinal}
Ist $S\subseteq T$ kofinal in $T$, so ist $\Hi(S)=\Hi(T)$.
Damit existiert $\sup S$ genau dann, wenn $\sup T$ existiert,
und die Werte stimmen überein.
Ist $S\subseteq T$ koinitial in $T$, so ist $\Lo(S)=\Lo(T)$;
die analoge Aussage gilt für die Infima.
\end{theorem}


\begin{theorem}[Hinzufügen kleinerer beziehungsweise größerer Elemente]
\label{b13guide:thm:one-sided-enlargement}
Für $S\subseteq P$ setze
\[
\begin{aligned}
 D(S)&=\{p\in P\mid\exists s\in S\;(p\leq s)\},\\
 A(S)&=\{p\in P\mid\exists s\in S\;(s\leq p)\}.
\end{aligned}
\]
Dann gilt
\[
 \Hi(D(S))=\Hi(S),\qquad \Lo(A(S))=\Lo(S).
\]
Somit ist $D(S)$ genau dann nach oben beschränkt wie $S$;
$A(S)$ ist genau dann nach unten beschränkt wie $S$.
Einschließlich der jeweiligen Existenzbedingungen gilt
\[
 \sup D(S)=\sup S,\qquad \inf A(S)=\inf S.
\]
Diese Aussagen schließen $S=\varnothing$ ein.
\end{theorem}

\begin{theorem}[Zwischenmengen und Hinzufügen vorhandener Grenzen]
\label{b13guide:thm:sandwich}
Sei $s=\sup S$ vorhanden. Für jede Menge $T$ mit
\[
 S\subseteq T\subseteq\mathord\downarrow s
\]
existiert $\sup T$, und $\sup T=s$.
Sei $m=\inf S$ vorhanden. Für jede Menge $T$ mit
\[
 S\subseteq T\subseteq\mathord\uparrow m
\]
existiert $\inf T$, und $\inf T=m$.
Insbesondere gilt im jeweiligen Existenzfall
\[
 \sup(S\cup\{\sup S\})=\sup S,\qquad
 \inf(S\cup\{\inf S\})=\inf S.
\]

Insbesondere gilt die \emph{Sandwichform}: Sind
$S\subseteq T\subseteq R$ und $\sup S=\sup R=s$,
so existiert $\sup T=s$. Sind $\inf S=\inf R=m$,
so existiert $\inf T=m$.

\end{theorem}

